\section{Results}
\label{sec:results}

\subsection{Evaluation populations and computational budgets}
One shared model was trained across 22 families with 402,336 presentations per family
(8,851,392 total), using 2,794 optimizer updates and deterministic FP32 execution on eight H100s.
Each family participated in 381 updates. The complete execution took 1,377.7 seconds, including
preparation and checkpoints (Table~\ref{tab:22-training}). The principal development population
contains 1,408 TRAIN-derived requests, 64 per family, evaluated using the original checkpoint and
the developed assembly, proposal and selection procedure. All requests, including failures,
remain in the denominator. This population was used during method development and is not an
independent held-out evaluation. Independent training-seed estimates and final held-out results
remain \PendingValue{independent-seed held-out evaluation}.

Separate paired diagnostics use 352 requests (16 per family), one draw per request and 64
sampling steps. They compare raw flow endpoints, terminal argmax and the bounded D1 completion
procedure, without the later full proposal and selection procedure used for the 1,408-request
population. These populations and inference budgets are not interchangeable. Means and sample
standard deviations across independent fits are not estimable from one original fit and its
continuations. Post-hoc rejection cannot increase a proposal's per-attempt acceptance probability
(Proposition~\ref{prop:posthoc-budget}); larger candidate pools and changed selection are therefore
reported separately from learned-model comparisons.

\subsection{Shared generation across 22 assembly families}
The selected development outputs achieve exact L1 in 1,387/1,408 requests (98.5\%), with all 22
families above 90\% (Table~\ref{tab:22-l1}). The lowest family yield is 58/64 (90.6\%) for aryl
reductive amination. Exact L1 requires molecular validation, an admitted decomposition and exact
forward replay. The equal-family and pooled yields coincide because every family has 64 requests.
These counts describe the combined model and inference procedure; the paired raw, null and cyclic
results for this same population remain unmeasured. Decomposition coverage, candidate precision,
ambiguity and the matched three-family comparison remain
\PendingValue{independent decomposition and common-program comparisons}. Acceptance checks alone
do not provide the precision of all attempted decompositions.

\begin{table}[H]
\caption{\textbf{Exact assembly across 22 ionizable-lipid reaction families.}
Selected exact-L1 percentages for one original checkpoint and 64 TRAIN-derived development
requests per family. Final includes the developed checked assembly, proposal and selection
procedure. Dashes denote unavailable paired measurements for this population, not zero.
The separate 352-request cyclic intervention is reported in Table~\ref{tab:22-controls}.
Independent-fit uncertainty is not estimated.}
\label{tab:22-l1}
\centering\footnotesize\setlength{\tabcolsep}{3pt}\renewcommand{\arraystretch}{1.12}
\begin{tabular}{@{}>{\raggedright\arraybackslash}p{.30\textwidth}ccccc@{}}
\toprule
Family & \shortstack{Raw\\L1 (\%)} & \shortstack{Final\\L1 (\%)} & \shortstack{Shared\\null (\%)} & \shortstack{Cyclic\\label (\%)} & \shortstack{Decomp. cov. /\\candidate prec.} \\
\midrule
A3 amine--aldehyde--alkyne & \PendingCell & 95.3 & \PendingCell & \PendingCell & \PendingCell \\
Acid--epoxide diester & \PendingCell & 98.4 & \PendingCell & \PendingCell & \PendingCell \\
AEMA & \PendingCell & 98.4 & \PendingCell & \PendingCell & \PendingCell \\
Aldehyde Ugi 3-CR & \PendingCell & 98.4 & \PendingCell & \PendingCell & \PendingCell \\
Aldehyde Ugi 4-CR & \PendingCell & 100.0 & \PendingCell & \PendingCell & \PendingCell \\
$\alpha$-Isocyanoester dihydroimidazole & \PendingCell & 100.0 & \PendingCell & \PendingCell & \PendingCell \\
Amine alkylation & \PendingCell & 96.9 & \PendingCell & \PendingCell & \PendingCell \\
Amine--epoxide & \PendingCell & 100.0 & \PendingCell & \PendingCell & \PendingCell \\
Aryl reductive amination & \PendingCell & 90.6 & \PendingCell & \PendingCell & \PendingCell \\
Aza-Michael acrylamide & \PendingCell & 100.0 & \PendingCell & \PendingCell & \PendingCell \\
Aza-Michael acrylate & \PendingCell & 100.0 & \PendingCell & \PendingCell & \PendingCell \\
Disulfide Michael & \PendingCell & 98.4 & \PendingCell & \PendingCell & \PendingCell \\
Epoxide O-acylation & \PendingCell & 98.4 & \PendingCell & \PendingCell & \PendingCell \\
iPhos & \PendingCell & 100.0 & \PendingCell & \PendingCell & \PendingCell \\
Ketone--isocyanide amide & \PendingCell & 100.0 & \PendingCell & \PendingCell & \PendingCell \\
Ketone Ugi 4-CR & \PendingCell & 96.9 & \PendingCell & \PendingCell & \PendingCell \\
Maleate & \PendingCell & 100.0 & \PendingCell & \PendingCell & \PendingCell \\
O-esterification & \PendingCell & 100.0 & \PendingCell & \PendingCell & \PendingCell \\
Passerini & \PendingCell & 100.0 & \PendingCell & \PendingCell & \PendingCell \\
Preassembled thiol-yne & \PendingCell & 100.0 & \PendingCell & \PendingCell & \PendingCell \\
Reductive amination & \PendingCell & 95.3 & \PendingCell & \PendingCell & \PendingCell \\
Thiolactone & \PendingCell & 100.0 & \PendingCell & \PendingCell & \PendingCell \\
\midrule
Equal-family macro & \PendingCell & 98.5 & \PendingCell & \PendingCell & \PendingCell \\
Pooled all requests & \PendingCell & 98.5 & \PendingCell & \PendingCell & \PendingCell \\
Worst-family result & \PendingCell & 90.6 & \PendingCell & \PendingCell & \PendingCell \\
\bottomrule
\end{tabular}
\end{table}

\subsection{Program conditioning and training diagnostics}
On the 352-request paired population, changing only the program identity at inference reduces D1
exact L1 from 140/352 to 105/352, a 9.94 percentage-point decrease. Raw endpoint exact counts are
80 versus 71, and terminal-argmax counts are 91 versus 85. The cyclic intervention retains the true
reaction core, roles, layout and depth. It therefore measures sensitivity to the program identity
within otherwise correct structural context, not the full effect of synthesis information.
The independently trained shared-null contrast remains \PendingValue{trained null comparison}.

Additional training did not improve D1 exact yield. With another 31,680 presentations per family,
220 large-batch updates and 2,640 small-batch updates produced 135 and 122 exact outputs,
respectively, versus 140 for their common initializer (all /352). Hardware, computation and
PCGrad grouping differed, so matched exposure does not imply matched cost. A second experiment
compared two 66-update continuations, each adding 9,504 presentations per family. Ordinary
categorical and coarse parent-group training gave 136 and 135 exact outputs, respectively.
Although coarse training reduced nonconsecutive-parent CAL errors from 1,164 to 1,026 of 7,936,
it did not preserve the original generation floors. Neither continuation replaced the original
checkpoint (Figure~\ref{fig:22-training}; Table~\ref{tab:22-architecture}).

\subsection{Decoder contributions and limitations}
For the original checkpoint on the 352-request diagnostic, raw endpoint, terminal-argmax and D1
readouts yield 80, 91 and 140 exact products, respectively. A fixed fallback among these already
computed candidates increases D1 exact yield to 145/352 and limited-design yield from 111 to
116/352, without new model or chemistry calls. This is a selection gain. In a separate paired
192-request repair diagnostic covering A3 and the two Ugi 4-CR families, exact L1 remains 187/192
while limited-design passes increase from 140 to 146. Additional attachment proposals change three
structures but add no design passes. These repairs were not substituted into the 1,408-request
routing population. Table~\ref{tab:22-decoder} retains the separate denominators and adverse results.
Matched whole-graph versus factorized capacity, cross-attention and gradient-control contrasts
remain \PendingValue{architecture comparisons}.

\subsection{Precursor novelty and diversity}
The selected population contains 1,407 distinct products among 1,408 requests. Product novelty
relative to TRAIN occurs in 1,216 observations (1,215 distinct identities). Exact products with at
least one globally TRAIN-novel recovered component occur in 1,050/1,408 requests (74.6\%),
representing 1,049 distinct products. This gives 745.0 distinct component-novel products per 1,000
requests. Role-relative novelty occurs in 1,081 observations (1,080 distinct products).
Component counts remain uneven: the pooled isocyanide role contains 29 identities among 190
observations, with Shannon and inverse-Simpson effective counts of 10.21 and 6.09
(Table~\ref{tab:22-diversity}). High product uniqueness therefore does not ensure broad component use.

TRAIN-relative novelty does not establish catalogue escape, recovery of deliberately withheld
components or source-disjoint generalization. Only 15/2,059 guarded CAL molecules are wholly
component-disjoint. The newer 1,859-molecule exact-source and representation-qualified CAL union
contains 12 such cases, all in two aza-Michael families; it differs from the older 1,821-reference
distribution cohort below. Final held-component and source-disjoint generation results remain
\PendingValue{generalization evaluation} (Table~\ref{tab:22-generalization}).

\subsection{Structural quality and agreement with lipid references}
All 1,408 selected structures are valid and connected, but only 1,324/1,408 (94.0\%) satisfy exact
L1 together with all applicable limited design checks (Table~\ref{tab:22-quality-metrics}).
There are 69 failures of the combined ring predicate. The scoped retained-head predicate fails in
5/64 aldehyde Ugi 4-CR requests; the other 1,344 structures are outside its qualified scope and
are not counted as passes. Chemistry-context review is flagged for 41/1,408 structures (2.9\%).
The overlapping flags include hydrated or hemiacetal carbons (24), peroxide groups (8), carbons
bearing at least three oxygens (6), other allenes (5), a small unsaturated ring (1) and a heteroatom
cumulene (1). Limited design without context flags holds for 1,292/1,408 (91.8\%).

\begin{table}[H]
\caption{\textbf{Structural quality of the selected development population.}
Counts retain all 1,408 requests except the explicitly scoped head predicate. The limited-design
conjunction includes exact L1 and the qualified chemical, head and ring checks; it is not a broad
empirical lipid-realism assessment. Dashes retain unavailable paired or separately resolved metrics.}
\label{tab:22-quality-metrics}\centering\footnotesize
\setlength{\tabcolsep}{3pt}\renewcommand{\arraystretch}{1.12}
\begin{tabular}{@{}>{\raggedright\arraybackslash}p{.61\textwidth}rc@{}}\toprule
Metric & Selected count & (\%)\\\midrule
Valid molecules / requests & 1,408/1,408 & 100.0 \\
Connected molecules / requests & 1,408/1,408 & 100.0 \\
Exact L1 + qualified limited design / requests & 1,324/1,408 & 94.0 \\
Limited design without context flags / requests & 1,292/1,408 & 91.8 \\
Justified narrow chemical alerts / requests & 0/1,408 & 0.0 \\
Chemistry-context flags / requests & 41/1,408 & 2.9 \\
Required head retained / scoped requests & 59/64 & 92.2 \\
Combined ring predicate / requests & 1,339/1,408 & 95.1 \\
Separate role-local cycle allocation & \PendingCell & \PendingCell \\
Separate requested ring-size agreement & \PendingCell & \PendingCell \\
TRAIN atom-environment coverage & \PendingCell & \PendingCell \\
Ring-system TRAIN support & \PendingCell & \PendingCell \\
Exact L1 + product/component features observed in TRAIN & 974/1,408 & 69.2 \\
\bottomrule\end{tabular}
\par\smallskip\raggedright\scriptsize
Source-positive controls pass the narrow chemical-alert screen in 29/29 cases, but cover only two families
and two studies, with zero individually characterized products in that control set. The other 20
families lack these controls. Absence of an alert does not establish stability, pKa or delivery activity.
Definitions and family-level counts: Tables~\ref{tab:22-structure}--\ref{tab:22-metric-definitions}.
\end{table}

The distribution comparison uses 1,821 exact-source-qualified \emph{synthetic CAL structures};
these are not an independent empirical lipid reference. Equal-family fingerprint precision is
56.9\%, coverage is 49.1\%, and mean within-family ECFP4 diversity is 0.552
(Table~\ref{tab:22-realism-metrics}). Precision varies from 1/64 for ketone Ugi 4-CR to 64/64 for
aza-Michael acrylate. Aryl reductive amination has 2/63 supported distinct outputs. These failures
remain visible despite high exact L1. A3 has 36/64 supported outputs; its normalized Wasserstein
distances for cLogP and rotatable bonds are 1.133 and 1.658. Ketone Ugi 4-CR has an sp$^3$-carbon
fraction distance of 1.082. All 24 descriptor panels are reported by family
(Table~\ref{tab:22-descriptors}). Independent empirical-reference comparisons, grouped two-sample
classification and blinded structural review remain unmeasured.

\begin{table}[H]
\caption{\textbf{Distributional agreement with synthetic CAL references.}
Values are equal-family means of 22 within-family statistics, not pooled molecular statistics or
independent-seed means. Each arm has 64 requests per family; fingerprint calculations deduplicate
connected products within family. The predecessor is the saved selection preceding the current
context-preserving selection, not raw model output. TRAIN controls are a reference diagnostic,
not independent quality evidence.}
\label{tab:22-realism-metrics}\centering\footnotesize
\setlength{\tabcolsep}{3pt}\renewcommand{\arraystretch}{1.12}
\begin{tabular}{@{}>{\raggedright\arraybackslash}p{.53\textwidth}ccc@{}}\toprule
Metric & Selected & Predecessor & TRAIN control\\\midrule
Fingerprint precision (\%) & 56.9 & 57.1 & 75.1 \\
Fingerprint coverage (\%) & 49.1 & 49.4 & 59.0 \\
Fingerprint recall (\%) & 82.1 & 81.4 & 88.4 \\
Distinct supported products / requests (\%) & 56.9 & 57.1 & 75.1 \\
Nearest-reference Tanimoto similarity & 0.638 & 0.641 & 0.693 \\
Within-family ECFP4 diversity & 0.552 & 0.550 & 0.529 \\
Descriptor distances (24 features, per family) & \multicolumn{3}{c}{Table~\ref{tab:22-descriptors}} \\
Independent descriptor precision / coverage & \PendingCell & \PendingCell & \PendingCell \\
Independent source-grouped classifier AUC & \PendingCell & \PendingCell & \PendingCell \\
\bottomrule\end{tabular}
\par\smallskip\raggedright\scriptsize
ECFP4 uses Morgan radius 2 and 2,048 bits; neighborhood metrics use five neighbors.
Similarity to this synthetic reference is not a probability of chemical realism.
Family counts, estimators and limitations: Tables~\ref{tab:22-realism} and~\ref{tab:22-metric-definitions}.
\end{table}

\subsection{Recursive computational makeability and strict dossiers}
The primary route criterion requires every assembly component to be directly vendor-listed or to
have a complete computational path to vendor-listed terminal materials. Planner-generated paths
and forward-applied published transforms may contribute without exact-substrate experimental
precedent. On the unchanged 1,408-product population, 616/1,408 (43.75\%) meet this criterion,
including 606 satisfying limited design and 604 also free of context flags
(Table~\ref{tab:22-routes}). L2 readiness is 1,189/1,408 (84.4\%); it does not require terminal
listing closure. Direct-only L3 covers 119/1,408 (8.5\%). Strict exact-source/current-item dossiers
remain a separate secondary outcome, closing 27/1,408 requests (1.9\%).

Makeability ranges from 0/64 for reductive amination and 2/64 for ketone Ugi 4-CR to 58/64 for
ketone--isocyanide amide. The complete denominator includes 3,723 required component branches.
The latest listing evidence is dated September 25, 2026 (local date), with a 30-day policy window.
A vendor-directory listing does not establish stock, price or an experimentally executed route.
The primary count rose from 581 to 616 through additional evidence, without changing molecules,
L1 checks or structural predicates. Missing listings and bounded route searches remain unresolved,
not evidence of unsynthesizability. Neither outcome establishes experimental synthesis or delivery.
